\chapter{Implementation Details}
\label{app:implementation}

\section{Shared Configuration}
\label{app:configs}

All experiments use Euclidean routing instances with solver-generated labels. Unless otherwise
stated, the dataset is split into 90\% training and 10\% validation with random seed 42. The
categorical diffusion process uses \(T=1000\) training diffusion steps with a linear beta
schedule from \(10^{-4}\) to \(0.02\). Validation and CADO rollouts use a cosine subsampling
schedule over the diffusion chain.

\begin{table}[htbp]
\centering
\caption{Shared data and diffusion settings.}
\label{tab:app-shared-config}
\resizebox{\textwidth}{!}{%
\begin{tabular}{lllll}
\toprule
Problem & Dataset file & Instance size & Label source & Diffusion setting \\
\midrule
TSP-50 & \texttt{tsp50\_128000\_concorde.txt} & 50 cities & Concorde & categorical, \(T=1000\) \\
CVRP-50 & \texttt{cvrp50-50\_128000\_pyvrp.txt} & 50 customers, 1 depot & PyVRP & categorical, \(T=1000\) \\
MDVRP-50 & \texttt{mdvrp10-10x5-5\_64000\_pyvrp.txt} & 50 customers, 2--5 depots & PyVRP & categorical, \(T=1000\) \\
\bottomrule
\end{tabular}
}
\end{table}

\section{Supervised DIFUSCO Runs}
\label{app:sl-configs}

The supervised experiments train the DIFUSCO-style denoising model with cross-entropy on
solver-generated binary labels. AdamW is used with cosine learning-rate decay, weight decay
\(10^{-4}\), dropout 0, and gradient clipping at norm 1.0. The TSP run uses the small model
because it was trained locally; the CVRP run uses the paper-size model on A100; and the MDVRP
run uses the small model for the assignment-level extension.

\begin{table}[htbp]
\centering
\caption{Supervised training configurations used for the reported baselines.}
\label{tab:app-sl-configs}
\resizebox{\textwidth}{!}{%
\begin{tabular}{llllllll}
\toprule
Run & Model & Platform & Epochs & Batch & LR & Eval subset & Validation decoding \\
\midrule
TSP-50 SL & \(h=128,L=6\) & Apple M4 Pro MPS & 50 & 16 & \(2\times10^{-4}\) & 500 & greedy + 2-opt, \(M_{\mathrm{eval}}=50\) \\
CVRP-50 SL & \(h=128,L=6\) & VESSL A100 & 50 & 64 & \(2\times10^{-4}\) & 500 & greedy + per-route 2-opt, \(M_{\mathrm{eval}}=50\) \\
MDVRP-50 SL & \(h=128,L=6\) & VESSL A100 & 50 & 64 & \(2\times10^{-4}\) & 500 & assignment decode, \(M_{\mathrm{eval}}=50\) \\
\bottomrule
\end{tabular}
}
\end{table}

\begin{table}[htbp]
\centering
\caption{Reported supervised validation outcomes.}
\label{tab:app-sl-results}
\resizebox{\textwidth}{!}{%
\begin{tabular}{lllll}
\toprule
Run & Best validation metric & Final validation metric & Wall-clock & Notes \\
\midrule
TSP-50 SL & 4.46\% gap & 4.46\% gap & 17 h 47 min & per-epoch validation curve not recorded due to logging issue \\
CVRP-50 SL & 25.96\% gap & 35.51\% gap & 2 h 9 min & best at epoch 10 with 2-opt validation \\
MDVRP-50 SL & 81.89\% accuracy & 55.75\% accuracy & 57 min 44 s & assignment-level metric; 0\% capacity violation \\
\bottomrule
\end{tabular}
}
\end{table}

\section{CADO Fine-Tuning Runs}
\label{app:cado-configs}

CADO fine-tuning starts from the corresponding supervised checkpoint. All reported CADO runs
use REINFORCE, AdamW with learning rate \(10^{-4}\), no weight decay, gradient clipping at norm
1.0, LoRA rank 2, and one fully trainable final graph layer plus the output head. Each epoch
uses 128 sampled training instances, accumulated into batches of 32, giving four optimizer
updates per epoch. Unless otherwise stated, the rollout uses \(M_{\mathrm{train}}=5\) and a
cosine schedule.

\begin{table}[htbp]
\centering
\caption{CADO fine-tuning configurations used for the reported experiments.}
\label{tab:app-cado-configs}
\resizebox{\textwidth}{!}{%
\begin{tabular}{lllllllll}
\toprule
Run & Start checkpoint & Reward & Epochs & Updates & \(M_{\mathrm{eval}}\) & Eval subset & Platform & Notes \\
\midrule
TSP-50 CADO & TSP-50 SL & SR & 1000 & 4000 & 50 & 200 & Apple M4 Pro MPS & reported TSP CADO run \\
CVRP-50 CADO & CVRP-50 SL & LCR & 1000 & 4000 & 25 & 100 & VESSL A100 & best checkpoint at epoch 470 \\
MDVRP-50 CADO & MDVRP-50 SL & SR & 500 & 2000 & 25 & 100 & VESSL A100 & assignment-level fine-tuning \\
\bottomrule
\end{tabular}
}
\end{table}

\begin{table}[htbp]
\centering
\caption{Reported CADO validation outcomes.}
\label{tab:app-cado-results}
\resizebox{\textwidth}{!}{%
\begin{tabular}{lllll}
\toprule
Run & Best validation metric & Final validation metric & Wall-clock & Main interpretation \\
\midrule
TSP-50 CADO & 3.77\% gap & 3.77\% gap & about 6--7 h & 1000-epoch run improves over 4.46\% SL baseline \\
CVRP-50 CADO & 20.24\% gap & 47.16\% gap & 7 h 1 min & best checkpoint improves, final checkpoint collapses \\
MDVRP-50 CADO & 80.86\% accuracy & 80.48\% accuracy & 1 h 49 min & recovers assignment accuracy with 0\% capacity violation \\
\bottomrule
\end{tabular}
}
\end{table}

\section{Evaluation Protocol}
\label{app:evaluation-protocol}

TSP and CVRP use the relative gap
\[
    \mathrm{Gap}(\%)
    =
    \frac{L_{\mathrm{pred}} - L_{\mathrm{gt}}}{L_{\mathrm{gt}}}
    \times 100.
\]
For TSP, \(L_{\mathrm{pred}}\) is computed after greedy heatmap decoding and 2-opt. For CVRP,
the supervised baseline reported in the main table uses greedy decoding with per-route 2-opt,
whereas the CADO run is evaluated with the implemented CADO CVRP decoder. CVRP validation also
records over-capacity route violations; all reported CVRP validation rows have 0\%
over-capacity rate. MDVRP reports assignment accuracy and capacity-violation rate rather than
route-cost gap, because full route construction within each depot is not part of the current
MDVRP experiment.

\section{Configuration Files and Logs}
\label{app:config-files}

The Hydra configuration files used by the code are stored under \texttt{configs/}. The main
entry-point configuration files are:
\begin{itemize}
    \item supervised: the three \texttt{difusco\_*\_config.yaml} files;
    \item CADO: the three \texttt{cado\_*\_config.yaml} files.
\end{itemize}
The exact run values reported above are taken from the experiment logs under \texttt{logs/},
because some runs override the YAML defaults from the command line or from the remote execution
environment.
